\chapter{Introducción}

Todos tenemos una idea intuitiva de qué es un objeto: es una cosa que podemos percibir, manipular y con la que podemos interactuar.

Por ejemplo, pensemos en un fibrón de pizarra. Podemos describir este objeto a partir de sus características: tiene una
determinada masa, color, forma, tamaño, etc. Estas características nos permiten describir el estado del objeto.

Pero un fibrón no solamente tiene características. También podemos realizar determinadas acciones con él. Podemos, por
ejemplo, escribir o dibujar sobre una pizarra utilizando el fibrón. También podemos realizar acciones sobre el propio
objeto, como ponerle o quitarle la tapa.

Un fibrón, además, no necesariamente tiene que ser considerado como un objeto aislado. Podemos observar que está
formado por distintas partes: un cuerpo, una tapa, una punta, etc. Cada una de estas partes podría, a su vez, ser
considerada como un objeto.

También podemos observar que un fibrón pertenece a una categoría más general de objetos: los útiles de escritura. Dentro
de esta categoría podemos encontrar lápices, lapiceras, marcadores, crayones, etc. Estos objetos comparten determinadas
características y comportamientos, aunque cada uno tenga sus particularidades.

Esta forma de pensar los objetos —qué características tienen, qué pueden hacer, cómo se relacionan con otros objetos y
qué tienen en común— es precisamente la base de la Programación Orientada a Objetos.

La definición de un objeto en el desarrollo de software no es muy diferente. Los objetos de software pueden no ser cosas
tangibles que podamos agarrar, percibir o tocar, pero son \textbf{modelos} de algo que puede realizar determinadas acciones y sobre el cual
se pueden realizar determinadas acciones. Formalmente, un objeto es una \textbf{colección de datos y comportamientos asociados}.

Veremos a lo largo de este apunte que los conceptos que utilizamos para describir el fibrón tienen su equivalente en POO.

\section{De la Programación Estructurada a la POO}
Hasta el momento, siempre tuvimos por un lado nuestros datos almacenados en variables, como enteros, listas, diccionarios, etc.
Y por otro lado, los procedimientos, descriptos en funciones, que reciben, procesan y devuelven un resultado. Esta forma de estructurar
el programa es conocido como programación estucturada.

Sin embargo, a medida que los problemas de software aumentan en complejidad, la separación entre
datos y funciones dificulta el mantenimiento y la escalabilidad del código.
La \textit{Programación Orientada a Objetos} (POO) propone una forma distinta de organizar y estructurar
el código: agrupar los datos y las acciones que los manipulan dentro de una misma entidad llamada \textit{objeto}.

\begin{figure}[!ht]
  \centering
  \begin{tikzpicture}[node distance=2cm, auto]
    % Estructurada
    \node [draw, rectangle, fill=gray!20, minimum width=3cm, minimum height=1cm] (datos) {Datos};
    \node [draw, rectangle, fill=gray!20, minimum width=3cm, minimum height=1cm, right=1.5cm of datos] (func) {Funciones};
    \draw[<->, thick] (datos) -- node[above] {\small Separados} (func);
    
    % POO
    \node [draw, rectangle, fill=blue!10, minimum width=4cm, minimum height=2.5cm, right=2.5cm of func] (objeto) {};
    \node [above] at (objeto.north) {\textit{Objeto}};
    \node [draw, rectangle, fill=blue!30, minimum width=3cm, baseline] at ($(objeto.center)+(0,0.5)$) {Estado (Datos)};
    \node [draw, rectangle, fill=blue!30, minimum width=3cm, baseline] at ($(objeto.center)+(0,-0.5)$) {Comportamiento};
  \end{tikzpicture}
  \caption{Comparación entre Programación Estructurada y Programación Orientada a Objetos.}
  \label{fig:estructurada_vs_poo}
\end{figure}

Como se ilustra en la \fig{fig:estructurada_vs_poo}, la POO permite integrar el estado y el comportamiento
en un bloque cohesivo, facilitando el modelado de sistemas complejos.

\section{¿Qué es un Objeto?}
En el mundo real, un objeto es cualquier entidad física o abstracta que posee características y realiza acciones.
Por ejemplo, un fibrón de pizarra tiene características como color, marca y masa, y realiza acciones como
escribir o ser borrado. Del mismo modo, un automóvil tiene un modelo, patente, kilometraje y color, y acciones
como acelerar, frenar o doblar.

En programación, un \textit{objeto} es un modelo conceptual que almacena información y define comportamientos
asociados. Todo objeto se compone de dos elementos fundamentales:

\begin{itemize}[leftmargin=*]
  \item \textit{Atributos}: representan el estado o las características del objeto (por ejemplo, el saldo de una cuenta bancaria).
  \item \textit{Métodos}: representan las acciones o procedimientos que el objeto puede ejecutar (por ejemplo, depositar o retirar dinero).
\end{itemize}

\section{Clases y Objetos: Moldes e Instancias}
Para crear objetos en un programa, primero debemos definir su estructura mediante una \textit{clase}.
Una clase funciona como un plano o plantilla arquitectónica que especifica qué atributos y métodos
tendrán los objetos construidos a partir de ella.

Un \textit{objeto} es una \textit{instancia} concreta de una clase. Aunque dos personas tengan un automóvil
del mismo modelo (construidos con el mismo plano), cada automóvil es un objeto independiente con su propia
patente, kilometraje acumulado y estado de desgaste.

\begin{figure}[!ht]
  \centering
  \begin{tikzpicture}[node distance=2.5cm]
    \node [draw, rectangle, fill=orange!20, minimum width=3.5cm, minimum height=1.5cm, align=center] (clase) 
      {\textbf{Clase (Plano)}\\ \texttt{Automovil}};
    
    \node [draw, rectangle, fill=green!20, minimum width=3.5cm, minimum height=1.2cm, align=center, below left=1cm and 0.5cm of clase] (obj1) 
      {\textit{Instancia 1}\\ Auto A (Verde)};
    
    \node [draw, rectangle, fill=green!20, minimum width=3.5cm, minimum height=1.2cm, align=center, below right=1cm and 0.5cm of clase] (obj2) 
      {\textit{Instancia 2}\\ Auto B (Rojo)};
    
    \draw[->, thick] (clase) -- node[left] {\small Instanciación} (obj1);
    \draw[->, thick] (clase) -- node[right] {\small Instanciación} (obj2);
  \end{tikzpicture}
  \caption{Relación entre una Clase (plantilla) y sus Instancias (objetos reales).}
  \label{fig:clase_objeto}
\end{figure}

Como se observa en la \fig{fig:clase_objeto}, una única definición de clase permite instanciar múltiples objetos
con valores de atributos independientes en memoria.
